% ---------------------------------------------------------------------------
\section{Reproducibility and Artefact Map}
\label{app:reproducibility}
% ---------------------------------------------------------------------------

The pipeline is deterministic. Every stage reads the previous stage's artefact and
writes a new one under \texttt{artifacts/}; no stage depends on notebook state or a
manual edit. The commands below are executed from the repository root, and the Python
environment is managed by \texttt{uv} so that the lock file, rather than a
hand-maintained list, defines the dependencies.

\begin{verbatim}
uv sync                                    # create the pinned environment
make pipeline                              # data, cleaning, features, models,
                                           # evaluation, monitoring, dictionary
uv run pytest -q                           # unit and integration tests
uv run ruff check src app scripts tests    # static checks
cd report && latexmk -pdf -interaction=nonstopmode -halt-on-error main.tex

# individual stages, when only one artefact needs regenerating
uv run python scripts/01_generate_context.py
uv run python scripts/02_clean_data.py
uv run python scripts/03_build_features.py
uv run python scripts/04_train_models.py
uv run python scripts/05_run_evaluation.py
uv run python scripts/06_run_monitoring.py
uv run python scripts/07_export_data_dictionary.py

# services (the console falls back to an in-process pipeline if the API is down)
uv run uvicorn revaluation.api.main:app --host 127.0.0.1 --port 8000
uv run streamlit run app/streamlit_app.py
\end{verbatim}

Table~\ref{tab:artefacts} maps the claims in this report to the artefact supporting
them. A reader who wants to check a number can open the corresponding file; a reader
who wants to re-derive it can re-run the stage that produces it.

\begin{table}[htbp]
\centering
\caption{Artefact map: the evidence behind each part of the report. Paths are
relative to the repository root.}
\label{tab:artefacts}
\small
\begin{tabularx}{\linewidth}{@{}L{0.46\linewidth}>{\raggedright\arraybackslash}X@{}}
\toprule
\textbf{Artefact} & \textbf{Content used in this report} \\
\midrule
\path{artifacts/metrics/model_comparison.csv} & Out-of-fold benchmark of the eight families (Table~\ref{tab:benchmark}) \\
\path{artifacts/metrics/training_summary.json} & Selected model, metrics, conformal quantiles, temporal holdout \\
\path{artifacts/metrics/evaluation_diagnostics.json} & Residual tests, SHAP drivers, subgroup performance \\
\path{artifacts/metrics/interval_coverage.csv} & Coverage and width by level and price band \\
\path{artifacts/metrics/cleaning_report.json} & Missing counts, imputation state, outlier rules \\
\path{artifacts/metrics/eda_summary.json} & Distributions, correlations, tiers, decades, VIF summary \\
\path{artifacts/metrics/shap_importance.csv} & Global attribution by raw feature and family \\
\path{artifacts/monitoring/monitoring_summary.json} & Reference vintage, window performance, drift, alerts \\
\path{artifacts/monitoring/alerts.json} & Triggered rules and the retraining decision \\
\path{artifacts/monitoring/monitoring_report.html} & Static monitoring dashboard \\
\path{docs/data_dictionary.md} & Field-level data dictionary (Appendix~\ref{app:dictionary}) \\
\path{artifacts/reports/ui_result.png} & Console screenshot (Figure~\ref{fig:console}) \\
\bottomrule
\end{tabularx}
\end{table}

Two practical notes complete the record. The base competition files are read from
their original directory and never modified, so the pipeline can be re-run against a
fresh download and will produce identical results. The generator and the model both
use a fixed seed of 42; changing it changes the simulated context but not the
structure of the analysis, and the monitoring thresholds and alert rules are
independent of it.
